\documentclass[11pt,a4paper]{article}
\usepackage{amsmath, amssymb, amsthm}
\usepackage{xcolor}
\usepackage{geometry}
\usepackage{booktabs}
\usepackage{array}
\usepackage{hyperref}
\usepackage{enumitem}
\usepackage{parskip}

\geometry{margin=2.5cm}

\hypersetup{
  colorlinks=true,
  linkcolor=blue!60!black,
  urlcolor=blue!60!black,
}

\newcolumntype{L}[1]{>{\raggedright\arraybackslash}p{#1}}

\title{%
  \textbf{Equation Reference} \\[0.4em]
  \large Optimal Execution of Portfolio Transactions \\
  \normalsize Almgren \& Chriss (December 2000)
}
\author{Symbol glossary and usage context for every numbered equation}
\date{}

\begin{document}
\maketitle
\tableofcontents
\bigskip

\noindent\textbf{Context.} This paper solves the problem of liquidating a large block of $X$ shares over a fixed horizon $T$ without moving the market too adversely. It balances two competing forces: selling slowly limits market impact (lower expected cost) but exposes you to price risk (higher variance). The result is a family of optimal trading schedules called the \emph{efficient frontier of execution strategies}.
% -----------------------------------------------------------------------
\section{Terms (\texorpdfstring{\S}{§}1)}
% -----------------------------------------------------------------------

\noindent\textbf{Drift.} Drift is when an event, for example a earnings report could have an impact on the underlying stock price, this event will make the stock either move up or down, in this case we need to take into account the drift and assume something when we use the formulas under 
\noindent\textbf{Bid-ask spread.} The difference between the price at which you can buy (ask) and sell (bid) a security. It represents a cost of trading and is an important consideration in execution strategies.
\noindent\textbf{Static vs dynamic strategy.} A static strategy is when before starting the trades it is already decided how the trades will be executed over the desired time. In contrast a dynamic strategy will adapt the way it liquidates the portfolio based on market conditions and price movements. 
\noindent\textbf{Naïve / VWAP strategy.} A naïve strategy, also known as a VWAP (Volume Weighted Average Price) strategy, executes trades evenly over the trading horizon, aiming to match the average market volume profile. It does not adapt to market conditions and is considered a static strategy.
\noindent\textbf{VaR.} Value at Risk (VaR) is a measure of the potential loss in value of a portfolio over a defined period for a given confidence interval. It is used to quantify the price risk associated with holding a position.
\noindent\textbf{L-VaR.} Liquidity-adjusted Value at Risk (L-VaR) extends the traditional VaR measure by incorporating the impact of liquidity on potential losses. It accounts for the fact that large trades can move the market, increasing the risk of adverse price movements.
\noindent\textbf{Risk Aversion.} Denoted by $\lambda$, it represents the trader's tolerance for risk. A higher value of $\lambda$ indicates greater aversion to price risk, leading to more conservative trading strategies that prioritize minimizing variance over expected cost.
% -----------------------------------------------------------------------
\section{Price Dynamics (\texorpdfstring{\S}{§}1)}
% -----------------------------------------------------------------------

\subsection*{Equation (1) — Actual Mid-Price Dynamics}

\begin{equation}
  S_k = S_{k-1} + \sigma\sqrt{\tau}\,\xi_k - \tau\,g\!\left(\frac{n_k}{\tau}\right)
\end{equation}

\textbf{When used.} Defines how the mid-price evolves at each discrete time step~$k$.
The three terms are: random diffusion, plus the previous price, minus the permanent
market impact from your own trading.  Used as the starting point for all cost expressions.

\begin{center}
\begin{tabular}{@{}L{3.5cm}L{10cm}@{}}
\toprule
\textbf{Symbol} & \textbf{Meaning} \\
\midrule
$S_k$ & Mid-price at the end of period $k$ \\
$S_{k-1}$ & Mid-price at the start of period $k$ \\
$\sigma$ & Annualised volatility (daily: $\sigma/\sqrt{252}$) \\
$\tau$ & Length of one interval; $\tau = T/N$ \\
$\xi_k$ & i.i.d.\ random shock, $\mathbb{E}[\xi_k]=0$, $\operatorname{Var}[\xi_k]=1$ \\
$n_k$ & Shares traded in interval $k$ (positive = sell) \\
$g(v)$ & Permanent market-impact function; $v = n_k/\tau$ is the trade rate \\
$\tau g(\cdots)$ & Permanent price depression per interval \\
\bottomrule
\end{tabular}
\end{center}

\subsection*{Equation (2) — Effective Execution Price}

\begin{equation}
  \tilde{S}_k = S_{k-1} - h\!\left(\frac{n_k}{\tau}\right)
\end{equation}

\textbf{When used.} The actual price received per share when selling $n_k$ shares in
interval $k$.  Unlike the mid-price $S_k$, this includes the additional \emph{temporary}
slippage $h(\cdot)$ paid because your own order moves the market against you within the
period.

\begin{center}
\begin{tabular}{@{}L{3.5cm}L{10cm}@{}}
\toprule
\textbf{Symbol} & \textbf{Meaning} \\
\midrule
$\tilde{S}_k$ & Effective execution price received in period $k$ \\
$h(v)$ & Temporary market-impact function; $v = n_k/\tau$ \\
$h(n_k/\tau)$ & Immediate price depression within the interval \\
\bottomrule
\end{tabular}
\end{center}

% -----------------------------------------------------------------------
\section{Trading Cost Measurement (\texorpdfstring{\S}{§}2)}
% -----------------------------------------------------------------------

\subsection*{Equation (3) — Total Proceeds from Liquidation}

\begin{equation}
  \sum_{k=1}^{N} n_k\,\tilde{S}_k
  = X S_0
  - \sum_{k=1}^{N} \tau\,x_k\,g\!\left(\frac{n_k}{\tau}\right)
  - \sum_{k=1}^{N} n_k\,h\!\left(\frac{n_k}{\tau}\right)
\end{equation}

\textbf{When used.} Decomposes total cash received into: (1)~the ``paper value'' $XS_0$
if everything were sold at the initial price, minus (2)~permanent impact losses, minus
(3)~temporary impact losses.  This accounting identity underpins all subsequent cost
expressions.

\begin{center}
\begin{tabular}{@{}L{3.5cm}L{10cm}@{}}
\toprule
\textbf{Symbol} & \textbf{Meaning} \\
\midrule
$X$ & Total shares to liquidate \\
$S_0$ & Initial stock price at time $0$ \\
$x_k$ & Shares still held at the \emph{start} of interval $k$ (inventory) \\
$\sum \tau x_k g(\cdots)$ & Cumulative permanent impact cost over the horizon \\
$\sum n_k h(\cdots)$ & Cumulative temporary impact cost over the horizon \\
\bottomrule
\end{tabular}
\end{center}

\subsection*{Equation (4) — Expected Shortfall $E(x)$}

\begin{equation}
  E(x) = \sum_{k=1}^{N} \tau\,x_k\,g\!\left(\frac{n_k}{\tau}\right)
        + \sum_{k=1}^{N} n_k\,h\!\left(\frac{n_k}{\tau}\right)
\end{equation}

\textbf{When used.} The expected total transaction cost (implementation shortfall) for a
given trading trajectory $x=\{x_1,\ldots,x_N\}$.  Random shocks cancel in expectation,
leaving only the two deterministic impact terms.  Minimising $E(x)$ alone gives the
constant-rate (VWAP) solution.  One of two objectives on the efficient frontier.

\begin{center}
\begin{tabular}{@{}L{3.5cm}L{10cm}@{}}
\toprule
\textbf{Symbol} & \textbf{Meaning} \\
\midrule
$E(x)$ & Expected cost of strategy $x$, relative to paper value \\
$\sum \tau x_k g(\cdots)$ & Permanent impact contribution \\
$\sum n_k h(\cdots)$ & Temporary impact contribution \\
\bottomrule
\end{tabular}
\end{center}

\subsection*{Equation (5) — Variance of Shortfall $V(x)$}

\begin{equation}
  V(x) = \sigma^2 \sum_{k=1}^{N} \tau\,x_k^2
\end{equation}

\textbf{When used.} The variance of total realised P\&L around expected cost.  Because
random shocks are i.i.d.\ and each remaining inventory $x_k$ is exposed to one period's
noise, variance accumulates as $\sigma^2\tau$ per share squared.  Minimising $V(x)$ alone
gives immediate liquidation.  The second objective on the efficient frontier.

\begin{center}
\begin{tabular}{@{}L{3.5cm}L{10cm}@{}}
\toprule
\textbf{Symbol} & \textbf{Meaning} \\
\midrule
$V(x)$ & Variance of cost for strategy $x$ \\
$\sigma^2$ & Per-period variance of price move \\
$\tau x_k^2$ & Variance contribution of interval $k$; proportional to inventory$^2$ \\
\bottomrule
\end{tabular}
\end{center}

% -----------------------------------------------------------------------
\section{Linear Impact Model (\texorpdfstring{\S}{§}3)}
% -----------------------------------------------------------------------

\subsection*{Equation (6) — Linear Permanent Impact}

\begin{equation}
  g(v) = \gamma v
\end{equation}

\textbf{When used.} Specifies permanent price depression proportional to trade rate
$v=n_k/\tau$.  The linearity makes the efficient frontier analytically tractable.
$\gamma$ is calibrated from empirical bid-ask and depth data.

\begin{center}
\begin{tabular}{@{}L{3.5cm}L{10cm}@{}}
\toprule
\textbf{Symbol} & \textbf{Meaning} \\
\midrule
$\gamma$ & Permanent impact coefficient (price/shares$^2$) \\
$v = n_k/\tau$ & Trade rate (shares per unit time) in interval $k$ \\
\bottomrule
\end{tabular}
\end{center}

\subsection*{Equation (7) — Linear Temporary Impact}

\begin{equation}
  h\!\left(\frac{n_k}{\tau}\right) = \varepsilon\operatorname{sgn}(n_k) + \frac{\eta}{\tau}\,n_k
\end{equation}

\textbf{When used.} Specifies the within-period execution cost.  The $\varepsilon$ term
is a fixed bid-ask half-spread regardless of size; the $\eta/\tau$ term is
size-proportional slippage.  Together they make trades expensive when done quickly.
Substituted into $E(x)$ to give Eq.~(8).

\begin{center}
\begin{tabular}{@{}L{3.5cm}L{10cm}@{}}
\toprule
\textbf{Symbol} & \textbf{Meaning} \\
\midrule
$\varepsilon$ & Fixed per-share cost (half bid-ask spread) \\
$\operatorname{sgn}(n_k)$ & $+1$ for sell, $-1$ for buy; ensures cost is always positive \\
$\eta$ & Temporary impact coefficient (price$\cdot$time/shares$^2$) \\
$\eta/\tau$ & Slippage per share; larger when $\tau$ is small (fast trading) \\
$\tilde{\eta}$ & Shorthand $\tilde{\eta} = \eta - \tfrac{1}{2}\gamma\tau$ used to simplify later formulas \\
\bottomrule
\end{tabular}
\end{center}

\subsection*{Equation (8) — Expected Cost, Linear Model}

\begin{equation}
  E(x) = \tfrac{1}{2}\gamma X^2
        + \varepsilon \sum_{k=1}^{N} |n_k|
        + \frac{\tilde{\eta}}{\tau} \sum_{k=1}^{N} n_k^2
\end{equation}

\textbf{When used.} Substituting the linear impact functions (Eqs.~6--7) into $E(x)$
(Eq.~4) yields this closed form.  The first term $\tfrac{1}{2}\gamma X^2$ is constant
across all strategies and cannot be optimised away.  The second and third terms depend on
the trade schedule $\{n_k\}$ and are targets for optimisation.

\begin{center}
\begin{tabular}{@{}L{3.5cm}L{10cm}@{}}
\toprule
\textbf{Symbol} & \textbf{Meaning} \\
\midrule
$\tfrac{1}{2}\gamma X^2$ & Unavoidable permanent impact; constant across all strategies \\
$\varepsilon\sum|n_k|$ & Total fixed-cost (spread) component; equals $\varepsilon X$ for a monotone sell \\
$(\tilde{\eta}/\tau)\sum n_k^2$ & Total size-proportional slippage; penalises large, fast trades \\
\bottomrule
\end{tabular}
\end{center}

% -----------------------------------------------------------------------
\section{Boundary Strategies \& Efficient Frontier (\texorpdfstring{\S}{§}4)}
% -----------------------------------------------------------------------

\subsection*{Equations (9--11) — Minimum-Impact (Constant-Rate) Trajectory}

\begin{align}
  n_k &= \frac{X}{N}, \qquad x_k = \frac{(N-k)X}{N} \\[4pt]
  E_{\mathrm{VWAP}} &= \tfrac{1}{2}\gamma X^2 + \varepsilon X
                      + \frac{\tilde{\eta} X^2}{T}\!\left(1+\frac{1}{N}\right) \\[4pt]
  V_{\mathrm{VWAP}} &= \tfrac{1}{3}\sigma^2 X^2 T
                      \!\left(1-\frac{1}{N}\right)\!\left(1-\frac{1}{2N}\right)
\end{align}

\textbf{When used.} Solution to $\min E(x)$ ignoring variance.  Sells $X/N$ shares in
each of the $N$ periods.  Minimises slippage by spreading trades evenly, but keeps large
inventory exposed to risk for a long time.  One extreme of the efficient frontier (minimum
cost, maximum variance).  As $N\to\infty$, $V_{\mathrm{VWAP}}\to\tfrac{1}{3}\sigma^2X^2T$.

\begin{center}
\begin{tabular}{@{}L{3.5cm}L{10cm}@{}}
\toprule
\textbf{Symbol} & \textbf{Meaning} \\
\midrule
$N$ & Total number of trading intervals \\
$T$ & Total liquidation horizon \\
$1/N$ terms & Finite-grid corrections; vanish as trading becomes continuous \\
$\tfrac{1}{3}\sigma^2X^2T$ & Limiting variance: scales with $\sigma^2$, position$^2$, and time \\
\bottomrule
\end{tabular}
\end{center}

\subsection*{Equations (12--13) — Minimum-Variance (Immediate) Liquidation}

\begin{align}
  n_1 &= X, \quad n_2 = n_3 = \cdots = n_N = 0 \\[4pt]
  E_{\mathrm{imm}} &= \varepsilon X + \frac{\eta X^2}{\tau}, \qquad V_{\mathrm{imm}} = 0
\end{align}

\textbf{When used.} Solution to $\min V(x)$: sell everything immediately.  Variance is
exactly zero (no time exposed to price moves) but execution cost is very high because the
entire position is dumped in one period, incurring maximum slippage $\eta/\tau$.  The
other extreme of the efficient frontier.

\begin{center}
\begin{tabular}{@{}L{3.5cm}L{10cm}@{}}
\toprule
\textbf{Symbol} & \textbf{Meaning} \\
\midrule
$\eta X^2/\tau$ & Cost of trading $X$ shares in one interval of length $\tau$ \\
$V=0$ & Zero risk because the position is closed immediately \\
\bottomrule
\end{tabular}
\end{center}

\subsection*{Equation (14) — Constrained Efficient-Frontier Problem}

\begin{equation}
  \min_x\; E(x) \quad \text{subject to}\quad V(x) \leq V^*
\end{equation}

\textbf{When used.} Formal statement of the efficient frontier problem.  For each risk
budget $V^*$, find the trajectory $x$ that minimises expected cost without exceeding that
variance.  Tracing out all $V^*$ from $0$ to $V_{\mathrm{VWAP}}$ sweeps the full
frontier.  Practitioners use this form when they have a hard risk constraint.

\subsection*{Equation (15) — Unconstrained (Lagrangian) Problem}

\begin{equation}
  \min_x\; U(x) = E(x) + \lambda\,V(x)
\end{equation}

\textbf{When used.} Equivalent reformulation of Eq.~(14).  The Lagrange multiplier
$\lambda$ converts the variance constraint into a penalty.  Each $\lambda\geq 0$
corresponds to one point on the frontier: $\lambda=0$ recovers the constant-rate
strategy; $\lambda\to\infty$ recovers immediate liquidation.  This is the workhorse
equation for deriving the optimal trajectory.

\begin{center}
\begin{tabular}{@{}L{3.5cm}L{10cm}@{}}
\toprule
\textbf{Symbol} & \textbf{Meaning} \\
\midrule
$\lambda$ & Risk-aversion parameter (Lagrange multiplier); $\lambda \geq 0$ \\
$U(x)$ & Composite objective (utility) to minimise \\
$\lambda V(x)$ & Variance penalty; scales with trader's risk aversion \\
\bottomrule
\end{tabular}
\end{center}

% -----------------------------------------------------------------------
\section{Optimal Trajectory Solution (\texorpdfstring{\S}{§}5)}
% -----------------------------------------------------------------------

\subsection*{Equation (16) — Optimality Condition (Euler--Lagrange)}

\begin{equation}
  \frac{x_{j-1} - 2x_j + x_{j+1}}{\tau^2} = \tilde{\kappa}^2\,x_j
\end{equation}

\textbf{When used.} First-order condition from differentiating $U(x)$ w.r.t.\ each $x_j$.
It is a second-order linear difference equation — the discrete analogue of $x''=\kappa^2 x$.
Its general solution involves hyperbolic functions.  Here
$\tilde{\kappa}^2 = \lambda\sigma^2/\tilde{\eta}$ encodes the trade-off between risk and
temporary impact cost.

\begin{center}
\begin{tabular}{@{}L{3.5cm}L{10cm}@{}}
\toprule
\textbf{Symbol} & \textbf{Meaning} \\
\midrule
$\tilde{\kappa}^2$ & $\tilde{\kappa}^2 = \lambda\sigma^2/\tilde{\eta}$; controls speed of liquidation \\
$(x_{j\pm1}-2x_j)/\tau^2$ & Discrete second derivative of the inventory trajectory \\
\bottomrule
\end{tabular}
\end{center}

\subsection*{Equation (17) — Optimal Inventory Trajectory \texorpdfstring{$\bigstar$}{}}

\begin{equation}
  x_j = X\,\frac{\sinh\!\bigl(\kappa(T-t_j)\bigr)}{\sinh(\kappa T)}
\end{equation}

\textbf{When used.} Solution to Eq.~(16) satisfying boundary conditions $x_0=X$,
$x_N=0$.  Tells you exactly how many shares to hold at each time $t_j = j\tau$.  The
shape is convexly declining: sell faster at the start (to cut variance) and slow down near
the end.  When $\kappa$ is large (high risk aversion) liquidation is nearly immediate;
when $\kappa$ is small the schedule approaches constant-rate.

\begin{center}
\begin{tabular}{@{}L{3.5cm}L{10cm}@{}}
\toprule
\textbf{Symbol} & \textbf{Meaning} \\
\midrule
$x_j$ & Shares held at time $t_j = j\tau$ \\
$\kappa$ & Continuous decay rate; $\kappa\approx\tilde{\kappa}$ for small $\tau$ \\
$\sinh(\cdot)$ & Hyperbolic sine; $\sinh u = (e^u - e^{-u})/2$ \\
$\sinh\!\bigl(\kappa(T-t_j)\bigr)$ & Remaining ``urgency'' at time $t_j$ \\
$\sinh(\kappa T)$ & Normalisation using total horizon \\
$\theta = 1/\kappa$ & Trade half-life: larger $\theta$ means slower execution \\
\bottomrule
\end{tabular}
\end{center}

\subsection*{Equation (18) — Optimal Trade List \texorpdfstring{$\bigstar$}{}}

\begin{equation}
  n_j = X\,\frac{2\sinh\!\bigl(\tfrac{1}{2}\kappa\tau\bigr)}{\sinh(\kappa T)}
        \cdot \cosh\!\Bigl(\kappa\bigl(T - t_{j-\frac{1}{2}}\bigr)\Bigr)
\end{equation}

\textbf{When used.} Derived from Eq.~(17) as $n_j = x_{j-1} - x_j$.  Gives the actual
order size to submit at each interval.  The $\cosh$ factor produces a U-shaped schedule:
larger trades at the start and end of the session.  This is the \emph{actionable output}
of the model — the schedule of order sizes to pass to the trading desk.

\begin{center}
\begin{tabular}{@{}L{3.5cm}L{10cm}@{}}
\toprule
\textbf{Symbol} & \textbf{Meaning} \\
\midrule
$n_j$ & Shares to trade in interval $j$ \\
$\cosh(\cdot)$ & Hyperbolic cosine; $\cosh u = (e^u + e^{-u})/2$ \\
$t_{j-\frac{1}{2}}$ & Midpoint of interval $j$; $t_{j-\frac{1}{2}} = (j-\tfrac{1}{2})\tau$ \\
$2\sinh(\tfrac{1}{2}\kappa\tau)/\sinh(\kappa T)$ & Normalisation ensuring $\sum_j n_j = X$ \\
\bottomrule
\end{tabular}
\end{center}

\subsection*{Equation (19) — Approximation for $\kappa$}

\begin{equation}
  \kappa \approx \sqrt{\frac{\lambda\sigma^2}{\tilde{\eta}}}
  \qquad\text{(continuous-time limit, small $\tau$)}
\end{equation}

\textbf{When used.} In the continuous limit $\tau\to 0$, $\tilde{\kappa}\to\kappa$ and
$\kappa$ takes this simple form.  It makes the model interpretable: $\kappa^2$ increases
with risk aversion $\lambda$ and volatility $\sigma^2$, decreases with temporary impact
cost $\tilde{\eta}$.  The trade half-life $\theta = 1/\kappa$.

\subsection*{Equation (20) — Efficient Frontier in Closed Form}

\begin{align}
  E(\kappa) &= \tfrac{1}{2}\gamma X^2
              + \tilde{\eta}\kappa X^2
                \frac{\coth(\kappa T)}{\tanh(\kappa\tau/2)}
                \cdot\tfrac{1}{2} \\[4pt]
  V(\kappa) &= \frac{\tfrac{1}{2}\sigma^2 X^2\,
                \bigl[T\coth(\kappa T) - \tau\coth(\kappa\tau/2)\bigr]}
               {\sinh^2(\kappa T)}
\end{align}

\textbf{When used.} Substituting Eqs.~(17)--(18) back into $E(x)$ and $V(x)$ gives these
formulas parameterised by $\kappa$ (equivalently $\lambda$).  As $\kappa$ traces from $0$
to $\infty$ the pair $(V(\kappa), E(\kappa))$ traces the entire efficient frontier in
(risk, cost) space.  Used to plot the frontier and find the optimal $\kappa$ for a given
risk tolerance.

\begin{center}
\begin{tabular}{@{}L{3.5cm}L{10cm}@{}}
\toprule
\textbf{Symbol} & \textbf{Meaning} \\
\midrule
$\coth(\cdot)$ & Hyperbolic cotangent; $\coth u = \cosh u/\sinh u$ \\
$\tanh(\cdot)$ & Hyperbolic tangent; $\tanh u = \sinh u/\cosh u$ \\
$\sinh^2(\kappa T)$ & Normalisation denominator in variance formula \\
\bottomrule
\end{tabular}
\end{center}

% -----------------------------------------------------------------------
\section{Alternative Risk Criteria (\texorpdfstring{\S}{§}6)}
% -----------------------------------------------------------------------

\subsection*{Equation (21) — Mean-Variance Utility}

\begin{equation}
  U_{\mathrm{util}}(x) = E(x) + \lambda_u\,V(x)
\end{equation}

\textbf{When used.} A mean-variance utility where the trader chooses $\lambda_u$ to
reflect their risk preferences.  Identical in form to Eq.~(15).  Natural choice for a
trader with quadratic utility; the frontier parameterised by $\lambda_u$ is exactly the
efficient frontier.

\subsection*{Equation (22) — Liquidity-Adjusted Value at Risk (L-VaR)}

\begin{equation}
  \mathrm{VaR}_p(x) = E(x) + \lambda_v\,\sqrt{V(x)}
\end{equation}

\textbf{When used.} An alternative risk criterion for risk managers.  $\lambda_v = z_\alpha$
is the $\alpha$-quantile of the standard normal (e.g.\ $\lambda_v\approx1.65$ for 95\%
confidence), so $\mathrm{VaR}_p$ is the $\alpha$-quantile of implementation shortfall.
Under L-VaR the optimal schedule differs from mean-variance because the penalty is on
$\sqrt{V}$ rather than $V$.

\begin{center}
\begin{tabular}{@{}L{3.5cm}L{10cm}@{}}
\toprule
\textbf{Symbol} & \textbf{Meaning} \\
\midrule
$\lambda_v$ & Normal quantile $z_\alpha$ (e.g.\ $1.65$ for 95\%, $2.33$ for 99\%) \\
$\sqrt{V(x)}$ & Standard deviation of cost \\
$E + \lambda_v\sqrt{V}$ & Mean plus scaled std dev $=$ upper confidence bound on cost \\
\bottomrule
\end{tabular}
\end{center}

% -----------------------------------------------------------------------
\section{Price Drift \& Serial Correlation (\texorpdfstring{\S}{§}7)}
% -----------------------------------------------------------------------

\subsection*{Equation (23) — Price Dynamics with Drift}

\begin{equation}
  S_k = S_{k-1} + \sigma\sqrt{\tau}\,\xi_k + \alpha\tau - \tau\,g\!\left(\frac{n_k}{\tau}\right)
\end{equation}

\textbf{When used.} Adds a deterministic drift $\alpha$ (expected return per unit time)
to Eq.~(1).  If $\alpha>0$ the stock is expected to rise, so a rational seller should
hold longer.  Starting point for the drift-adjusted optimal strategy.

\begin{center}
\begin{tabular}{@{}L{3.5cm}L{10cm}@{}}
\toprule
\textbf{Symbol} & \textbf{Meaning} \\
\midrule
$\alpha$ & Expected price drift per unit time (may be negative) \\
$\alpha\tau$ & Expected price change per interval due to drift \\
\bottomrule
\end{tabular}
\end{center}

\subsection*{Equation (24) — Expected Cost with Drift}

\begin{equation}
  E(x) = \tfrac{1}{2}\gamma X^2
        - \alpha\sum_{k=1}^{N}\tau x_k
        + \varepsilon\sum_{k=1}^{N}|n_k|
        + \frac{\tilde{\eta}}{\tau}\sum_{k=1}^{N}n_k^2
\end{equation}

\textbf{When used.} Extension of Eq.~(8) to include drift.  The new term
$-\alpha\sum\tau x_k$ rewards holding inventory when $\alpha>0$ (drift reduces cost) and
penalises holding when $\alpha<0$, shifting the optimal trajectory accordingly.

\begin{center}
\begin{tabular}{@{}L{3.5cm}L{10cm}@{}}
\toprule
\textbf{Symbol} & \textbf{Meaning} \\
\midrule
$-\alpha\sum\tau x_k$ & Drift benefit/cost; positive drift reduces expected cost \\
$\sum\tau x_k$ & ``Time in market'' weighted by inventory; increases with slow execution \\
\bottomrule
\end{tabular}
\end{center}

\subsection*{Equation (25) — Optimal Static Holding with Drift}

\begin{equation}
  \bar{x} = \frac{\alpha}{2\lambda\sigma^2}
\end{equation}

\textbf{When used.} The inventory level at which the marginal benefit of holding (drift
$\alpha$) equals the marginal risk cost ($\lambda\sigma^2$).  In the drift-adjusted
framework the optimal trajectory converges toward $\bar{x}$ rather than zero.  If
$\bar{x}\geq X$, it is never optimal to sell at all.

\subsection*{Equations (26--27) — Optimal Trajectory and Trades with Drift}

\begin{align}
  x_j(\alpha) &= x_j(0) + \frac{\bar{x}}{X}\,X\Bigl[1 - \tfrac{x_j(0)}{X}\Bigr] + \cdots \\[4pt]
  n_j(\alpha) &= n_j(0) - \text{correction proportional to }\bar{x}
\end{align}

\textbf{When used.} Modifies the baseline trajectory (Eq.~17) and trade list (Eq.~18) to
account for the drift correction proportional to $\bar{x}/X$.  When drift is zero,
$\bar{x}=0$ and the baseline is recovered.  Used when the trader has a view on near-term
price direction.

\subsection*{Equation (28) — Serial Correlation Gain}

\begin{equation}
  \text{max gain per period} = \frac{\rho^2\sigma^2\tau^2}{4\eta}
\end{equation}

\textbf{When used.} When price returns have autocorrelation coefficient $\rho$, an
optimal strategy can exploit the predictability.  This gives the maximum extra expected
gain per interval from adjusting the trade schedule.  It vanishes when $\rho=0$ (i.i.d.\
returns) or when $\eta$ is very large (exploiting autocorrelation is too costly).

\begin{center}
\begin{tabular}{@{}L{3.5cm}L{10cm}@{}}
\toprule
\textbf{Symbol} & \textbf{Meaning} \\
\midrule
$\rho$ & Autocorrelation coefficient of period returns; $|\rho|\leq1$ \\
$\rho^2\sigma^2\tau^2$ & Predictable component of squared price move \\
$4\eta$ & Temporary-impact cost of exploiting the predictability \\
\bottomrule
\end{tabular}
\end{center}

% -----------------------------------------------------------------------
\section{Multi-Asset Portfolio Extension (\texorpdfstring{\S}{§}8)}
% -----------------------------------------------------------------------

\subsection*{Equation (29) — Multi-Asset Expected Cost}

\begin{equation}
  E[x] = \sum_{j=1}^{M}\left[
    \tfrac{1}{2}\gamma_j X_j^2
    + \varepsilon_j\sum_{k}|n_{jk}|
    + \frac{\tilde{\eta}_j}{\tau}\sum_{k}n_{jk}^2
  \right]
\end{equation}

\textbf{When used.} Extends Eq.~(8) to a portfolio of $M$ securities.  In the base
model, expected cost decomposes asset-by-asset with no cross-impact terms.  Each security
$j$ has its own parameters $\gamma_j$, $\varepsilon_j$, $\eta_j$.  Cross-asset coupling
enters only through the variance (covariance matrix).

\begin{center}
\begin{tabular}{@{}L{3.5cm}L{10cm}@{}}
\toprule
\textbf{Symbol} & \textbf{Meaning} \\
\midrule
$j$ & Index over securities $j = 1,\ldots,M$ \\
$X_j$ & Total shares of security $j$ to liquidate \\
$n_{jk}$ & Shares of security $j$ traded in period $k$ \\
$\gamma_j,\varepsilon_j,\eta_j$ & Impact parameters specific to security $j$ \\
\bottomrule
\end{tabular}
\end{center}

\subsection*{Equation (30) — Multi-Asset Variance \texorpdfstring{$\bigstar$}{}}

\begin{equation}
  V[x] = \sum_{k=1}^{N} \tau\,\mathbf{x}_k^\top\,C\,\mathbf{x}_k
\end{equation}

\textbf{When used.} The variance of total P\&L for a multi-asset portfolio, where $C$ is
the $M\times M$ covariance matrix of asset returns per unit time.  Off-diagonal elements
of $C$ capture diversification: correlated assets that are both sold amplify risk;
negatively correlated assets partially offset.  The single-asset formula Eq.~(5) is
recovered when $M=1$ and $C=\sigma^2$.

\begin{center}
\begin{tabular}{@{}L{3.5cm}L{10cm}@{}}
\toprule
\textbf{Symbol} & \textbf{Meaning} \\
\midrule
$\mathbf{x}_k$ & $M$-dimensional inventory vector at time $k$ (one entry per asset) \\
$C$ & $M\times M$ return covariance matrix (per unit time) \\
$\mathbf{x}_k^\top C\,\mathbf{x}_k$ & Portfolio variance in period $k$; a quadratic form \\
$\sum_k\tau(\cdots)$ & Cumulates variance contribution over all $N$ intervals \\
\bottomrule
\end{tabular}
\end{center}

\subsection*{Equation (31) — Multi-Asset Optimality Condition}

\begin{equation}
  H\,\frac{\mathbf{x}_{j-1} - 2\mathbf{x}_j + \mathbf{x}_{j+1}}{\tau^2}
  = \lambda C\,\mathbf{x}_j
\end{equation}

\textbf{When used.} Matrix generalisation of Eq.~(16).  $H$ is a diagonal matrix of
$\tilde{\eta}_j$ values; $C$ is the covariance matrix.  The solution requires
diagonalising $H^{-1}\lambda C$ (a generalised eigenvalue problem).  Each eigenvector of
$H^{-1}C$ defines a ``principal portfolio'' with its own $\kappa_i$ and optimal
trajectory; the full solution is a linear combination.

\begin{center}
\begin{tabular}{@{}L{3.5cm}L{10cm}@{}}
\toprule
\textbf{Symbol} & \textbf{Meaning} \\
\midrule
$H$ & Diagonal $M\times M$ matrix of temporary-impact coefficients $\tilde{\eta}_j$ \\
$C$ & $M\times M$ covariance matrix of returns \\
$H^{-1}C$ & Coupling matrix whose eigenvalues give the $\kappa_i$ values \\
$\lambda$ & Scalar risk-aversion parameter (same for the whole portfolio) \\
\bottomrule
\end{tabular}
\end{center}

\subsection*{Equation (32) — Diagonal Special Case}

\begin{equation}
  \Gamma_{jj} = \gamma_j, \qquad H_{jj} = \eta_j
\end{equation}

\textbf{When used.} When there is no cross-asset market impact, both impact matrices are
diagonal.  Eq.~(31) then reduces to $M$ independent single-asset problems, each solved by
its own $\sinh$-type trajectory.  Cross-asset effects enter only through $C$, coupling the
optimal speeds for correlated assets.

\begin{center}
\begin{tabular}{@{}L{3.5cm}L{10cm}@{}}
\toprule
\textbf{Symbol} & \textbf{Meaning} \\
\midrule
$\Gamma$ & $M\times M$ permanent impact matrix; diagonal $=$ no cross-asset impact \\
$H$ & $M\times M$ temporary impact matrix; diagonal $=$ no cross-asset impact \\
$\gamma_j, \eta_j$ & Per-asset impact scalars on the main diagonal \\
\bottomrule
\end{tabular}
\end{center}

\bigskip
\noindent\rule{\linewidth}{0.4pt}\\[0.3em]
\small Almgren \& Chriss, ``Optimal Execution of Portfolio Transactions,'' December 2000.
Equation reference compiled September 2026.

\end{document}